\section{Exact determinant transport of theta functions and their integral charts}
\label{sec:cubic-theta-transport}

The source's Weil twisting rule permits a stronger comparison than
a comparison of heights of CM models.  It gives an exact transport
between two automorphic character types.  The distinction between
transporting the type and leaving a projector fixed is essential.

\begin{theorem}[AC-G42: determinant transport in the source theta integral]
\label{AC-G42}
Let $E/F$ be a CM quadratic extension, and let $\nu$ be a
finite Hecke character of $E$, of order dividing $3$, satisfying
$\nu(\overline a)=\nu(a)^{-1}$.  For a hermitian space $V$,
write $G_V=U(V)$.  There is a finite automorphic character
\[
 \vartheta_V(g)=\nu(x),\qquad
       x/\overline x=\det g,                             \tag{G42.1}
\]
with values in $\mu_3$.  It is trivial on $G_V(F)$ and on
the archimedean factor.  Write $M_Vf=\vartheta_V f$.

Use the source's \texttt{WeilTied.om\_tw} and
\texttt{WeilTied.th} conventions, or its concrete
\texttt{WeilFun} realization.  For an abstract
\texttt{ThetaX} model also use its
\texttt{ThetaSmoothnessCondition}, so its evaluation
range is the smooth automorphic function space.
Let $V,W$ satisfy the source's \texttt{SigOK} condition
and have unequal ranks, let $(\chi,\chi')$ be an admissible splitting pair,
and put
$d=1$ if $\dim V<\dim W$, and $d=-1$ otherwise.
For every Schwartz vector $\Phi$ the kernels satisfy
\[
 K_{\chi\nu,\chi'\nu}(g,h;\Phi)
       =\vartheta_V(g)^{-d}\vartheta_W(h)^d
                       K_{\chi,\chi'}(g,h;\Phi).        \tag{G42.2}
\]
Accordingly, with exactly the same measure, Petersson
normalization, and Schwartz vector,
\[
 \Theta_{\chi\nu,\chi'\nu,\Phi}(M_V^d f)
          =M_W^d\Theta_{\chi,\chi',\Phi}(f).             \tag{G42.3}
\]
The equality applies whenever the source theta integral is
defined.  Its integrability is unchanged by this transport.
Multiplication by $M_V$ preserves cuspidality and all
archimedean differential operators and is an isometry for
the Petersson norm.

For a finite chain of increasing ranks, with both splitting
characters on every edge multiplied by $\nu$, the complete
theta composite satisfies
\[
        \Theta^{\nu}_{b,a}=M_b\Theta_{b,a}M_a^{-1}.
                                                               \tag{G42.4}
\]
For a chain of decreasing ranks the corresponding formula
uses $M_b^{-1}$ and $M_a$.  Each comparison uses the same
Schwartz vector on the corresponding old and new edge.
\end{theorem}
\begin{proof}
\proofstep{1. The determinant character and its normalization.}
On $\mathbb A_F^\times$ the anti-invariance relation gives
$\nu^2=1$, whereas $\nu^3=1$; hence $\nu=1$ there.
Adelic Hilbert 90 gives a root $x$ in (G42.1).  Here it can
be checked directly.  Locally, a norm-one element $z\ne-1$
has the root $1+z$, because $1+z=z\overline{(1+z)}$.
For $z=-1$, any nonzero element with conjugate equal to
its negative is a root.  In the split algebra the root of
$(a,a^{-1})$ can be chosen as $(a,1)$.  At almost all
places the extension is unramified and $z$ is a unit;
multiply a chosen root by an element of the base field
to make its valuation zero.  The local roots then form
an idele.  The same formulas over the number field give
a rational root for every rational norm-one element.

Two roots differ by an element of $\mathbb A_F^\times$,
on which $\nu$ is trivial.  Thus (G42.1) is independent
of the root and is multiplicative.  The local formulas,
translated from a neighborhood of the identity, show
continuity; away from finitely many places one can use
unit roots and $\nu$ is unramified.  Rational roots give
triviality on $G_V(F)$.  At infinity the unitary groups
are connected, so a continuous character with finite
image is trivial.  This is exactly the source convention
\texttt{IsDetRoot}, namely $x=(\det g)\overline x$.

\proofstep{2. The actual Weil and theta formulas.}
The source's upward rules multiply the Weil operator
by $\vartheta_V(g)^{-1}$ when $\chi'$ is replaced by
$\chi'\nu$, and by $\vartheta_W(h)$ when $\chi$ is
replaced by $\chi\nu$.  The downward rules have the
opposite exponents.  Since $\nu$ is trivial on
$\mathbb A_F^\times$, the restriction condition
\texttt{ResEta} is unchanged, so the two rules may be
applied consecutively.  Evaluation on $\Phi$ at zero
gives (G42.2).

In the source convention the theta integral contains
$K(g,h;\Phi)f(g)$, with no complex conjugation of $f$.
Multiplying $f$ by $\vartheta_V(g)^d$ cancels the first
factor in (G42.2); the second is independent of the
integration variable.  Pulling it outside gives
(G42.3), with the same normalization constant.
All factors have absolute value one.  Thus the old and
new integrands have the same absolute value, and their
integrability conditions agree.

The determinant is one on every unipotent subgroup.
Every constant term of $M_V f$ is therefore
$\vartheta_V$ times the corresponding constant term
of $f$; cuspidality is preserved.  The character is
trivial at infinity, so archimedean derivatives commute
with $M_V$.  Its absolute value one proves the
Petersson isometry.  In the abstract source model,
\texttt{ThetaTied.eChr\_range} represents $\nu$ in
\texttt{X.Chr}, since its infinity exponents are zero.
Multiplication by $\vartheta_V$ preserves each clause
of \texttt{IsAutFunSm}: automorphy and a finite level,
smoothness, finite-dimensional center-stable spans
by conjugating such a span, and growth of all
archimedean derivatives.  Thus the assumed evaluation
range and evaluation injectivity give the same
operator $M_V$ on \texttt{X.AF}; the cusp equivalence
gives it on \texttt{X.cusp} as well.
Finally (G42.3) on consecutive
edges gives (G42.4) by cancellation of the intermediate
$M_j^{-1}M_j$.  The decreasing case is identical with
the exponents reversed.
\end{proof}

\begin{proposition}[AC-G42-lattice: fixed-level character types and chart isometries]
\label{AC-G42-lattice}
Fix a compact open subgroup $K\subset G_V(\mathbb A_{F,f})$
and a finite type $\tau$ on $K$, with the conventional
right-equivariance rule $f(gk)=\tau(k)^{-1}f(g)$.
Multiplication by $M_V$
maps automorphic forms of type $\tau$ to forms of type
$\tau\otimes\vartheta_V^{-1}|_K$.  It uses the same double-coset
representatives and the same rational stabilizer groups.
After adjoining $\mu_3$ to the coefficient field, the map
is an integral isomorphism at every finite coefficient
place for lattices specified by the source's pointwise
polynomial or jet-coordinate integrality test.  It is
an isometry at every infinite coefficient embedding
for the corresponding \texttt{NormTieG} expression.

More explicitly, if $g_c$ are the fixed representatives,
then the chart coordinate polynomial $P_c$ becomes
\[
                  P_c^{\nu}=\vartheta_V(g_c)P_c.         \tag{G42.5}
\]
This preserves both rational and integral polynomial
spans.  It preserves all lattice indices after any
finite extension of the coefficient DVR.  This
conclusion concerns the transported character type;
it does not assert that the character stays trivial
on the old $K$.

If one instead insists on trivial finite type, the
single subgroup $K\cap\ker\vartheta_V$ has index at
most $3$ in $K$.  It need not be a product of local
subgroups.  The product of the local kernels can have
index $3^s$ when $s$ local factors have nontrivial cubic
character; these are different level choices.
\end{proposition}
\begin{proof}
With a right-equivariance convention
$f(gk)=\tau(k)^{-1}f(g)$, multiplication by
$\vartheta_V(g)$ changes the right type to
$\tau\otimes\vartheta_V^{-1}|_K$.  Equivalently, the
right-equivariance multiplier is multiplied by
$\vartheta_V(k)$.  No subgroup or coset representative
has changed.

For a rational stabilizer element $\gamma$ satisfying
$\gamma g_c=g_c k$, character multiplicativity gives
$\vartheta_V(k)=\vartheta_V(\gamma)=1$.  Thus the
stabilizer relations on each polynomial fiber are
the old ones.  Evaluating the new form at $g_c$ gives
(G42.5).  Each multiplier is a root of unity in the
coefficient field, hence is a unit at every finite
place.  The old polynomial lies in the rational or
integral span if and only if the new one does.  The
same statement holds after a local basis change and
for jets, since the multiplier is constant in the
archimedean variable.  This is precisely the form of
the source's \texttt{ValuationCompatibleChart.rat}
and \texttt{.lat} conditions.

At every complex coefficient embedding the multiplier
has absolute value one.  The class weights and the
archimedean frame determinants in \texttt{NormTieG}
are unchanged, and each polynomial norm in its sum
is unchanged.  Hence the complete norm is unchanged.
An integral linear isomorphism identifies the quotient
modules defining every finite lattice index; it
continues to do so after every finite DVR extension.

Finally $K/(K\cap\ker\vartheta_V)$ is the image of
$K$ in $\mu_3$, and so has order at most $3$.  If
$K=\prod_v K_v$, the quotient by $\prod_v\ker
(\vartheta_V|_{K_v})$ is the product of the local
images.  It has order $3^s$ when exactly $s$ of those
images have order $3$.  This proves the level
distinction without replacing one subgroup by the other.
\end{proof}

\begin{proposition}[AC-G42-Hecke: the coefficient twist and the projector obligation]
\label{AC-G42-Hecke}
At a split finite place where $\nu$ and the level are
unramified, a degree-$j$ standard Hecke operator has
its eigenvalue multiplied under $M_V$ by
$\nu_P(\varpi)^j$.  Thus the characteristic-polynomial
coefficients acquire the powers prescribed in AC-G39,
with the Hecke character chosen in the same Artin
reciprocity convention as the desired Galois character.

Let $R$ be a coefficient DVR containing $\omega$,
and let $e$ be any rational finite Hecke projector on a
finite-dimensional space with full integral lattice
$L$, and let $A\subset L\otimes\operatorname{Frac}(R)$
be a full theta lattice.  Under an integral transport
$M$ as above, put
\[
 L^{\nu}=ML,\quad A^{\nu}=MA,\quad
                  e^{\nu}=MeM^{-1}.                    \tag{G42.6}
\]
Then $M(eL)=e^{\nu}L^{\nu}$ and
$M(eA)=e^{\nu}A^{\nu}$.  All signed indices of the
projected lattices are equal, also after finite DVR
extension.  Keeping $e$ fixed can destroy full rank;
there is an explicit three-dimensional example.
\end{proposition}
\begin{proof}
Write a Hecke operator as a sum of right translations
over representatives $h_a$ of a double coset.  Since
$\vartheta_V$ is trivial on the local compact subgroup,
it is constant on that double coset.  Directly,
\[
 \sum_a(M_Vf)(g h_a)
       =\vartheta_V(h_a)\vartheta_V(g)
                                    \sum_a f(g h_a).
\]
For the degree-$j$ diagonal double coset its determinant
is $\varpi^j$.  The split determinant root is the idele
with component $\varpi^j$ at $P$ and component $1$ at
its conjugate prime.  Formula (G42.1) gives the factor
$\nu_P(\varpi)^j$.  Any fixed scalar normalization of
the Hecke operator commutes with this calculation.

The two assertions following (G42.6) follow by applying
$e^{\nu}M=Me$.  The isomorphism $M$ identifies the
intersections and quotient modules used to define
signed indices, so their lengths and normalized
lengths agree over every finite DVR extension.

For the promised example, let $R$ contain $\omega$,
let $C_3=\langle s\rangle$, and use the function lattice
$R^{C_3}$.  Let $f_0(s^a)=1$, $f_1(s^a)=\omega^a$,
and let $M$ be multiplication by $f_1$.  It is an
integral isometry, and $Mf_0=f_1$.  Over the fraction
field, the projector
\[
              e_0=\frac{1+s+s^2}{3}
\]
for the translation action satisfies $e_0f_0=f_0$
and $e_0f_1=0$.  Thus the old selected line is killed
after the twist if its projector is not transported.
This is a projector-fragment counterexample, not a
countermodel to the complete source \texttt{NData}.
\end{proof}

\paragraph{Verified source locators and the exact remaining comparison.}
The supplied
\path{Formalism/LocalToGlobal/BesselHermitianTheta/TiedWeil.lean}
defines \texttt{IsDetRoot}, the four fields of
\texttt{TwRule}, \texttt{WeilTied.om\_tw}, and the
theta integral \texttt{WeilTied.th}; these are the
conventions used in (G42.1)--(G42.3).  The concrete
\path{Development/AbelianVarietyTheory/Definitions/AutomorphicMembership/AdelicDataRowMeasureScaleTrivialTuple.lean}
has the identical \texttt{WeilCore.Rules.twW/twV}
and \texttt{WeilFun.twW/twV} rules.  The separate
\texttt{SplitFacTwistRight.lean} proves the corresponding
split-factor identities.  These named declarations
are verified at the pinned commit used in AC-G41.
The source's \texttt{MinfDensity.lean},
\texttt{ThetaTied.eChr\_range}, supplies the character
realization, while \texttt{TiedSmooth.lean},
\texttt{ThetaSmoothnessCondition}, supplies the function
range.  The concrete \texttt{Realises.forms} field
in \texttt{AdelicDataRowMeasureScaleTrivialTuple.lean}
has precisely that latter range property.

The chart assertions use the actual
\path{Formalism/LocalToGlobal/HermitianForms/TiedTwGen.lean},
\texttt{ValuationCompatibleChart.rat/lat} and
\texttt{NormTieG}.  The determinant phases also appear
explicitly, with opposite exponents, in
\path{Development/AbelianVarietyTheory/Definitions/AutomorphicMembership/ReferenceCharactersFactorTame.lean},
\texttt{rhoAt}, \texttt{rhoAt'}, \texttt{xiDet\_eq},
and \texttt{rhoAt\_eq\_hecke/rhoAt'\_eq\_hecke}.
Its \texttt{KDq\_le\_ker\_of\_tame} permits a tame
character supported in the indicated finite set;
that particular factorization statement does not
require absolute residue degree one.

These results give an exact comparison for the theta
functions and character-type chart lattices constructed
by the transport.  To identify them with a prescribed
source family one must still identify its chosen
Schwartz vectors, finite type/chart, and projector with
the transported ones.  This is materially more specific
than a new CM-model existence problem.  It is not the
same-ambient comparison in
\texttt{IntegralHeckeForms.IdxAt}: that predicate keeps
its specified \texttt{eG}, \texttt{V0}, and theta data.
Finally \texttt{LatAtLattice.lean},
\texttt{GlobalTheory.Mem.Col2.TC.Later}, lists
\texttt{N : NStage B C W} as additional data; its
\texttt{tw} construction does not derive that field
from the theta transport.  The relative character
and Selmer tie between the separate packets therefore
remains an explicit gate.

\paragraph{Dictionary.}
The finite character retains the already audited
SGA~1, Expose~V, \S\S4 and~7 interpretation.
The unitary determinant character and its automorphic
function-space transport are explicit representations;
no additional geometric square, motive, or positivity
principle is inferred from their norm equality.
